\section{Related Work}

\paragraph{LLMs as proxies for human populations.}
A growing body of work uses language models conditioned on demographic or attitudinal profiles
to approximate human samples---``silicon sampling''~\citep{argyleSiliconSampling}---and surveys
the promise and limits of LLMs as tools for computational social science~\citep{ziems2024css}.
The evidence, however, is mixed. On the supportive side, profile-conditioned models can reproduce
aggregate survey \emph{distributions}~\citep{argyleSiliconSampling}, and agents grounded in a
two-hour interview \emph{per person} replicate that individual's survey responses at $\approx$85\% of
the person's own two-week test--retest reliability~\citep{park2024agents}---though both successes are
on \emph{stated attitudes}, and the latter buys its fidelity with rich per-person grounding rather
than the cheap demographic profile we test. On the cautionary side,
LLM ``survey responses'' latch onto spurious option ordering rather than
meaning~\citep{dominguezolmedo2024questioning}; persona variables explain less than 10\% of the
variance in subjective annotation and help only when the persona--outcome correlation is
strong~\citep{hu2024persona}; assigning personas can surface biases and \emph{degrade}
performance on objective tasks~\citep{zheng2024helpful}; and persona simulations tend toward
\emph{caricature}~\citep{cheng2023compost}; and ad-hoc persona \emph{generation} itself injects
systematic bias, with election and survey forecasts deviating from real
outcomes~\citep{li2025personacatch}. Our work contributes a direct sim-to-real
\emph{validity} test against measured real-world behaviour (not survey distributions), and—unlike
prior work focused on accuracy or annotation—isolates the marginal effect of persona conditioning
with a no-persona baseline on a copy-ranking task.

\paragraph{The Upworthy archive and headline performance.}
The Upworthy Research Archive~\citep{matiasUpworthy} has supported large-scale analyses of what
makes headlines succeed; for example, negativity in headlines increases consumption at
scale~\citep{robertson2023negativity}. These works model \emph{aggregate} drivers of clicks; we
instead use the archive as a held-out benchmark to test whether an external persona-based
simulator can \emph{predict} the within-test winner. Most directly, concurrent work on the same
archive~\citep{ye2024lola} reports that pure-LLM headline prediction (prompting, embeddings,
fine-tuning) only marginally beats random; our results refine this picture by showing that
\emph{after} filtering to A/B tests with a reliable winner, a no-persona LLM ranker is in fact a
medium-strength predictor---and that adding personas erases this. The finding that personas can
introduce harmful bias echoes work on LLM simulations producing
\emph{caricatures}~\citep{cheng2023compost}.

\paragraph{Construct validity and effect sizes.}
We frame the intent-vs-engagement mismatch using the classical construct-validity
framework~\citep{cronbach1955construct} and pre-register effect-size thresholds following
Cohen~\citep{cohen1988power}.
